\documentclass[11pt]{article}
\usepackage[T1]{fontenc}
\usepackage{textcomp}
\usepackage[margin=0.75in]{geometry}
\usepackage{amsmath,amssymb,amsthm}
\usepackage{array,booktabs}
\usepackage{hyperref}
\setlength{\parindent}{0pt}
\newtheorem{theorem}{Theorem}
\theoremstyle{definition}
\newtheorem{definition}{Definition}
\title{Flow Proof Artifact: Group-inv-right}
\author{Generated by \texttt{flow doc proof}}
\date{Flow Proof Book}
\begin{document}
\maketitle
Every group element has a right inverse.\\[0.5em]
\textbf{Source.} dummit-foote — *Abstract Algebra*, §1.1\\[0.5em]
\bigskip
\noindent{\large \textbf{Derived fact 1.} \textit{g * inv(g) = 1} \hfill Group \cdot inverse \cdot right inverse recovers the identity}\par
\smallskip\noindent\textit{Coordinate: Group \cdot inverse \cdot right inverse recovers the identity}\par
\smallskip\noindent\textit{Source: dummit-foote}\par
\smallskip\noindent\textit{Built on: left inverse recovers the identity, for inverse on Group, one is the left identity, for identity on Group, parentheses do not matter, for multiplication on Group}\par
\medskip
\noindent\textbf{Goal.} g * inv(g) = 1\par
\noindent$\displaystyle \forall g \in Group\quad g \cdot inv(g) = 1$\par
\medskip
\noindent
\begin{tabular}{@{} >{\raggedright\arraybackslash}p{0.06\textwidth} >{\raggedright\arraybackslash}p{0.47\textwidth} >{\raggedleft\arraybackslash}p{0.41\textwidth} @{}}
\textbf{\#} & \textbf{Proof} & \textbf{Mathematics} \\
\midrule
\textbf{1.} & We prove that right inverse recovers the identity for inverse on Group. & \\[0.45em]
\textbf{2.} & We invoke the definitional clause governing inverse on Group: left inverse recovers the identity, for inverse on Group (instantiated for g). & \\[0.45em]
\textbf{3.} & We invoke the definitional clause governing identity on Group: one is the left identity, for identity on Group (instantiated for inv(g)). & \\[0.45em]
\textbf{4.} & We invoke the definitional clause governing multiplication on Group: parentheses do not matter, for multiplication on Group (instantiated for g, inv(g), 1). & \\[0.45em]
\textbf{5.} & From step 2, step 3, and step 4, this implies g times inv(g) equals 1. Hence proven. & $\displaystyle g \cdot inv(g) = 1$ \\[0.45em]
\bottomrule
\end{tabular}
\label{thm:Group-inverse-right_inverse_recovers_the_identity}
\par\medskip\hrule
\end{document}
